\documentclass[12pt]{article}
\usepackage[utf8]{inputenc}
\usepackage[T1]{fontenc}
\usepackage{hyperref}
\title{Azhe Language Resources for K-12 Curriculum Development}
\author{Compiled by Claude Sonnet 5.5 (claude-sonnet-5-5) for Andr\'as Kornai}
\date{}
\begin{document}
\maketitle

\section{Introduction}
Azhe (ISO 639-3 yiz; Glottocode azhe1235) is a Loloish (Lolo-Burmese, Sino-Tibetan) language spoken by Yi people in Yunnan, China, in Mile, Huaning, Kaiyuan and Jianshui counties. It is not the same as Ache (yif). Speaker estimates vary from about 54,000 to 100,000. This report rests on one search round and is short because little turned up.

\section{Primary Educational Resources}
No primers, textbooks or school programmes for Azhe appeared in the search results.

\section{Linguistic and Reference Materials}
\begin{itemize}
\item Overview article, which cites Wang Chengyou (2003) for three mutually intelligible dialects: \url{https://en.wikipedia.org/wiki/Azhe_language}
\item Disambiguation from the related Ache Yi language: \url{https://en.wikipedia.org/wiki/Ache_Yi_language}
\end{itemize}

\section{Religious and Community Resources}
\begin{itemize}
\item A missionary-oriented people-group profile (PDF). It is not a language resource, and I did not check its content: \url{https://people-groups.asiaharvest.org/people-group-profiles/yunnan/azhe.pdf}
\end{itemize}
The Wikipedia article says Azhe once had a traditional script known as Baimawen, or the shaman's script.

\section{Digital and Audio Resources}
None found.

\section{Curriculum Development Context}
No evidence of Azhe-language teaching was found. Any curriculum would have to start from documentation work, using a standard Yi script or Chinese as the medium. This was not investigated.
\end{document}